\section{Introducción}
El presente informe desarrolla la Actividad 5 de Sistemas de Radiocomunicación, dedicada al radar \cite{consigna}. En la parte teórica se definen el radar y sus aplicaciones, la ecuación del radar y los conceptos de detección, la sección radar (RCS), el filtro apareado, el clutter y la señal chirp con su compresión, sobre la base del material de la cátedra \cite{apunte} y de la bibliografía indicada \cite{skolnik, richards, cumming}.

En la parte práctica se procesa la señal \senal{} asignada al equipo: una detección simulada de radar de 100 pulsos con entre 2 y 10 blancos de rangos y velocidades desconocidos. Se implementaron la compresión con filtro apareado, la ventana CFAR, los canceladores MTI simple, STI simple y MTI doble y el filtrado Doppler, y con ellos se determinaron el rango y la velocidad de cada blanco.
